\newcommand{\F}{\mathbb{F}}
\newcommand{\C}{\mathbb{C}}
\newcommand{\an}{\mathscr{A}} 
\newcommand{\Q}{\mathbb{Q}}
\renewcommand{\hom}{\operatorname{Hom}}
\newcommand{\ehom}{\operatorname{End}}
\colorlet{gris}{black!70}

\-\vspace{-0.3cm}
  {\hrule \vspace{-0.15cm}
  \centering \textbf{2025-II} \vspace{0.1cm}
  \hrule}

1. Defina ou enuncie em detalhe cada um dos ítens: \\
(a) Aplicação diff num ponto \(\x\);  %\textcolor{gris}{\(F: U\ab \subseteq \R^n \to \R^m\) é \emph{diff} em \(p \in \R^n\) sse \(\exists L \in \hom(\R^n, \R^m): F(p+v) = F(p) + Lv + o\|v\|\), quando \(v\to 0\). }\\ 
(b) \(\blacksquare\) (\emph{TVM\(-\R^n\)}); %\textcolor{gris}{Sejam \(f: U\ab \subseteq \R^n \to \R\) continua em \([p, p+v]\) e diff em \((p,p+tv)\), então \(\exists \theta \in (0,1)\) tal que \( f(p+v) - f(p) = df(p + \theta v) \cdot v \).   }\\ 
(c) \(\blacksquare\) (\emph{Função Inversa}); %\textcolor{gris}{Sejam \(U\ab \subseteq \R^n\), \(F \in C^k(U,\R^n)\) e \(p\in U\). Então, \(^kF: \V_p \ \leftrightharpoons \ DF(p)\) é isomorfismo.  }\\ 
(d) Submersão;  %\textcolor{gris}{Uma aplicação \(F\) diff em \(p\) é submersão (em \(p\)) sse \(DF(p): \R^n \twoheadrightarrow \R^m\). }\\
(e) \(\blacksquare\) (\emph{FLS});  % \textcolor{gris}{Seja \(F\in C^k(U,\R^m)\) submersão em \(p\in U\). Então, \(\exists \V_p \subseteq U\) e \(G:V\ab \simeq \V_p\) tais que \(F\circ G: V \ni (\x, x_{m+1}, \ldots, x_n) \mapsto \x \in \R^m\). }\\
(f) \(\blacksquare\) (\emph{Função Implícita}).\vspace{0.3cm}\\ %\textcolor{gris}{Sejam \(U\ab = V\times W \subseteq \R^n \times \R^m\), \(G \in C^k(U,\R^m)\) e \((p,q) \in U\). Se \(G(p,q)=0\) e \(\det JG_q(p,q) \neq 0\), então \(\exists \mathcal{W}_q\) e \(\exists F \in C^k(\V_p, \R^m)\) tais que \(\forall (\x,\y) \in \V_p\times \mathcal{W}_q, \ G(\x,\y) = G(\x, F(\x)) = 0\). }\\ 
2. Sejam \(F,G: U\ab \subseteq \R^n \to \R^n\) diff em \(p \in U\), tais que \(F(p) = G(p)\). Mostre que,
\[DF(p) = DG(p) \ \leftrightharpoons \ \lim_{v\to 0} \frac{F(p+v) - G(p+v)}{\|v\|} = 0. \]
\textcolor{gris}{
    Segue essencialmente da definição, pois 
    \[DF(p) - DG(p) = \lim_{v\to 0} \frac{F(p+v) -DG(p+v)}{\|v\|}\] 
}
\hspace{-0.1cm}3. Sejam \(U:= \{(x,y) \in \R^2 : x,y >0 \}\), \(p,q >1\) tais que \(\nicefrac{1}{p} + \nicefrac{1}{q} = 1\) e \(f: U \to \R\) tal que, 
\[f(x,y) := \frac{x^p}{p} + \frac{y^q}{q}.\] 
(a) Explique por qué \(\forall t>0\), a restrição de \(f\) à hipérbole \(xy = t\) possui um ponto mínimo global \((x_t, y_t)\). \\ 
\textcolor{gris}{O conjunto \(\{(x,y) \in \R^2 : xy=t\}=: H_t\ce\subseteq \R^2 \) e ilimitado, mas \(f\) é contínua e \(f(x,y) \to \infty\) quando \(\|(x,y)\| \to \infty\), por tanto \(f\) tem mínimo global \((x_t, y_t) \in U\). } \vspace{0.3cm}\\
(b) Utilizando \(\blacksquare\) (\emph{Mult. de Lagrange}), mostre que \(x_t^p = y_t^q\). \\ 
\textcolor{gris}{Seja \(g(x,y) = xy - t\), então \(\exists \lambda \in \R \) tal que,  %: df(x_t, y_t) = \lambda dg(x_t, y_t)\) e \(g(x_t,y_t) = 0\). 
\[ \begin{cases}
        x_t^{p-1}dx + y_t^{q-1}dy = \lambda (y_t dx + x_tdy) \\ 
        x_ty_t  - t = 0
    \end{cases}  \]
Isso \(\leftrightharpoons \ x_t^p = \lambda y_tx_t = \lambda x_ty_t = y_t^q = \lambda t  \).
} \vspace{0.3cm} \\ 
(c) Mostre que \(f(x_t, y_t) = t \) e conclua que \(\forall x,y > 0\), 
\[ xy \leq \frac{x^p}{p} + \frac{y^q}{q}.  \]
\- \vspace{-0.5cm}
\textcolor{gris}{
\begin{align*}
f(x_t, y_t) &=  \frac{x_t^p}{p} + \frac{y_t^q}{q} = x_t^p \left(\frac{1}{p} + \frac{1}{q}\right) = x_t^p\\ 
&= (x_t^p)^{\nicefrac{1}{p} + \nicefrac{1}{q}} = x_t (x_t^{p})^{\nicefrac{1}{q}} = x_ty_t = t. 
\end{align*}
Dados \(x,y > 0\), \(xy = t = f(x_t, y_t) \leq f(x,y) = \frac{x^p}{p} + \frac{y^q}{q}\). 
}\vspace{0.3cm}\\ 
4. Utilize \(\blacksquare\) (\emph{Função Implícita}) para mostrar que o sistema, 
\[\begin{cases}
    \textcolor{gray}{G_1(\x) = } \  w^2 + 2x^2+y^2 - z^2 - 6 = 0 \\ 
    \textcolor{gray}{G_1(\x) = } \  wxy - xyz = 0 
\end{cases}, \]
pode ser resolvido em termos de \(w = w(y,z)\) e \(x = x(y,z)\) numa vizinhança do ponto \((p,q) = ((y,z),(x,w)) = ((2,1),(1,1))\). Calcule as derivadas parciais de \(w\) e \(x\) em \(p\). \\
\textcolor{gris}{
    Sejam \(G(\x,\y) = (G_1(\x,\y), G_2(\x,\y)): \R^4 \to \R^2\). Note que \(G \in C^\infty (\R^4, \R^2) \), \(G(p,q) = (0,0)\) e,
    \[\det (JG_q(p,q)) = \det \begin{bmatrix}
        4 & 2 \\ 
        0 & 2
    \end{bmatrix} = 8 \neq 0.\]  
    Pelo \(\blacksquare\) (\emph{Função Implícita}) \(\exists \mathcal{W}_q\) e \(F\in C^\infty(\V_p, \R^2)\) onde \(\forall (\x, \y) \in \V_p \times \mathcal{W}_q, \ G(\x,\y ) = G(\x, F(\x))\), seja lá
    \[F(y,z) = (x(y,z), w(y,z)).\]  
    Pelo {\scriptsize \(\boxtimes\)} 8.5. \(JG_q(p,q)\cdot JF_p(p,q) = - JG_p(p,q)\), ou seja,
    \[ \begin{bmatrix}
        x_y(p,q) & x_z(p,q) \\
        w_y(p,q) & w_z(p,q) 
    \end{bmatrix}  = \begin{bmatrix}
        4 & 2 \\ 
        0 & 2
    \end{bmatrix}^{-1} \begin{bmatrix}
        -4 & 2 \\ 
        0 & 2
    \end{bmatrix} = \begin{bmatrix}
        -1 & 0 \\ 
        0& 1
    \end{bmatrix}.\]
}
5. Mostre que \(\omega \in \Omega^1(\R^2\setminus \{0\}) \) é fechada mas não exata,  
\[\omega := -\frac{y}{x^2 + y^2}dx + \frac{x}{x^2+y^2} dy \textcolor{gray}{ \ = fdx + gdy}.\] 
\textcolor{gris}{
    Lembre-se que \(\omega \) é \underline{fechada} \(\leftrightharpoons\ d\omega = 0 \), 
    \[ d\omega = (g_x - f_y) dx \wedge dy = \frac{2[x^2+y^2] - 2x^2 - 2y^2}{[x^2+y^2]^2}dx \wedge dy = 0. \] 
    Seja \(\E^1 \subseteq \R^2\), variedade compacta, orientada, sem borda, suponha por contradição que \(\omega \) for \underline{exata}, usando \(\psi: [0,2\pi) \ni t \mapsto (\cos(t), \sin(t)) \in \E^1\) temos,% então \(\int_{\E^1} \omega =0\), pelo {\scriptsize \(\boxtimes\)} (\emph{Stokes}) mas,  
    \[\int_{\E^1}w = \int_0^{2\pi} \sin^2(t) + \cos^2(t)  dt = \int_0^{2\pi}  dt = 2\pi \neq 0. \ \lightning\]
}
\hspace{-0.15cm}De um exemplo de \(_d\overline{M}\subseteq \R^n\) compacta que possui uma forma de volume exata. \\ 
\textcolor{gris}{
    Seja \(_2\overline{M} = \overline{\B^2} := \{\x \in \R^2: \|\x\|\leq 1\}\). Note que \(\omega_\text{vol} = dx \wedge dy\) e \(\omega = d(-\frac{y}{2}dx+\frac{x}{2}dy)\). 
}\vspace{0.3cm}\\
Porém, se \(M\) for compacta e \(\partial M = \emptyset \), então nenhuma forma de volume é exata. \\
\textcolor{gris}{
    Supõe por contradição que \(\omega_\text{vol} = d\eta \), então do \(\blacksquare\) (\emph{Stokes}), 
    \[0< \int_M\omega_\text{vol} = \int_{\partial M} \eta = 0. \ \lightning\] 
} 
\newpage

\-\vspace{-0.3cm} 
  {\hrule \vspace{-0.15cm}
  \centering \textbf{2025-I} \vspace{0.1cm}
  \hrule}

1. Sejam \(^\infty\gamma, F, G, H : \R^2 \simeq \R^2\) tais que \(\gamma = H \circ G \circ F\), \(F(0,0)= (1,1), \ G(1,1) = (2,2)\), e \(H(2,2) = (0,0)\). Determine \(DH^{-1}(0,0)\), sabendo que: 
{\small\[D\gamma(\textbf{0}) = \begin{bmatrix}
    2 & 1 \\ 
    1 &1 
\end{bmatrix}; \ DF(\textbf{0}) = \begin{bmatrix}
    3 & 3 \\ 
    2 &1 
\end{bmatrix}; \ DG(1,1) = \begin{bmatrix}
    \nicefrac{1}{2} & 0 \\ 
    0 & 2 
\end{bmatrix}\]  }
\par\noindent 
\textcolor{gris}{Temos inversas bem definidas, portanto,%lembrando que se \(F(p) = q\), então \([DF^{-1}(q)] = [DF(p)]^{-1}\), temos 
\[\gamma \circ F^{-1} \circ G^{-1} = H \ \leftrightharpoons\ H^{-1} = G \circ F\circ \gamma^{-1}.\]
Pelo \(\blacksquare\) (\emph{RC}) \(DH^{-1}(\textbf{0})= DG(F(\gamma^{-1}(\textbf{0}))) \cdot DF(\gamma^{-1}(\textbf{0})) \cdot D\gamma^{-1}(\textbf{0}) = DG(1,1)\cdot DF(\textbf{0}) \cdot [D\gamma(\textbf{0})]^{-1}\), 
{\small\[DH^{-1}(\textbf{0}) = \begin{bmatrix}
    \nicefrac{1}{2} & 0 \\ 
    0 & 2
\end{bmatrix}\begin{bmatrix}
    3 & 3 \\ 
    2 &1 
\end{bmatrix}\begin{bmatrix}
     2 & 1 \\ 
    1 &1 
\end{bmatrix}^{-1} = \begin{bmatrix}
0 & \nicefrac{3}{2} \\
2 & 0 
\end{bmatrix}.\]  }
} \\ 
2. Seja \(F\in C^1(\R^3, \R^3) \) tal que \(\forall \x \in \R^3,\ DF(\x)\in \operatorname{GL}_3(\R)\) e \(F^{-1}(K)\) é compacto sempre \(K\subset \R^3 \) for compacto. Mostre que \(^1F: \R^3 \simeq \R^3\).  \\ 
\textcolor{gris}{
Segue do \(\blacksquare\) (\emph{Função Inversa}) que \(F\) é uma aplicação aberta, em particular, \((\operatorname{Im}F)\ab \subset \R^3\). Seja \(\y \in \overline{\operatorname{Im}F}\), então \(\exists (F(\x_n)) \subset \operatorname{Im}F : F(\x_n) \to \y \). Considere,
\[K := \{F(\x_n): n \in \N\} \cup \{\y\}\subset \R^3.\] \(K\) é compacto pois pra qualquer \((\textbf{k}_n)\subset K\), temos \(\textbf{k}_n \to \textbf{k} \in K\) e seqüência convegente é limitada. Agora, \((\x_n) \subset F^{-1}(K)\) compacto por hipotese, logo \(\exists \x_{\phi(k)} \to \x \in F^{-1}(K)\ce\). Por continuidade \(F(\x_{\phi(k)}) \to F(\x) = \y\) (toda subseqüência converge ao mesmo límite). Finalmente, \(\operatorname{Im}F\) é tanto aberto quanto fechado em \(\R^3\) conexo \(\leftrightharpoons\) \(\operatorname{Im}F = \R^3\). \vspace{-0.1cm}\\ } 
\hspace{-0.2cm}3. (a) Seja \(F: U\ab \subset \R^n \to \R^m\) uma aplicação. Defina a derivada de \(F\) no punto \(\x\in U\).  \\ 
\textcolor{gris}{\(F \) é diff em \(\x\) sse \(\exists L \in \hom(\R^n, \R^m) : F(\x+v) = F(\x) + Lv + o\|v\|\), quando \(v\to 0\), cujo caso \(DF(\x) = L = JF(\x)\). } \vspace{-0.1cm} \\ 
\hspace{-0.1cm}(b) Seja \(B :\R^n \times \R^m \to \R^p\) bilinear. Mostre que \(B\) é diff e exiba a derivada de \(B\).  \\ 
\textcolor{gris}{Note que \(B((\x,\y) +(h,k))= B(\x + h, \y +k) = B(\x,\y) + B(\x, k) + B(h, \y) + B(h,k)\),
\begin{align*}
\lim_{\|(h,k)\| \to 0} \frac{B(h,k)}{\|(h,k)\|_1} &\leq \lim_{\|(h,k)\| \to 0} \frac{C\|h\|_1\|k\|_1}{\|(h,k)\|_1}\\ 
&\leq \lim_{\|(h,k)\| \to 0} C\|k\|_1 = 0.  
\end{align*} 
Logo, \(B\) é diff e \(DB(\x,\y)(h,k) = B(\x, k) + B(h,\y)\), linear em \((h,k)\) de fato. }  \vspace{0.3cm}\\ 
(c) Seja \(F : \R^{n^2}\cong M_n(\R) \ni X \mapsto XX^t \in \operatorname{Sym}_n(\R) \cong \R^{\nicefrac{n(n+1)}{2}}\). Mostre que \(\id_n\) é valor regular de \(F\). 
\textcolor{gris}{
\begin{align*}
 F = \mu \circ G: M_n(\R)  &\overset{G}{\longrightarrow} M_n(\R) \times M_n(\R) \overset{\mu}{\longrightarrow} M_n(\R) \\ 
 X\hspace{0.5cm} &\mapsto \hspace{1cm}(X,X^t) \hspace{0.8cm}\mapsto \hspace{0.3cm}XX^t 
\end{align*}
Note que \(G\) é linear quanto \(\mu\) é bilinear, do \(\blacksquare\) (\emph{Regla da Cadeia}) e ítem (b), 
\begin{align*} DF(X)(H) &= D\mu(G(X))DG(X)(H) \\ 
    &= D\mu(X,X^t)(H,H^t) = XH^t + HX^t.\end{align*}
Sejam \(X \in M_n(\R) : F(X) = XX^t = \id_n\) e \(S \in \operatorname{Sym}_n(\R)\) e \(H:= \nicefrac{1}{2}SX\), então  
\[DF(X)(H) = \nicefrac{1}{2}X(SX)^t + \nicefrac{1}{2}(SX)X^t = S. \]
Portanto \(DF(X)\) é sobrejetiva e \(\id_n\) valor regular. 
}  \vspace{0.3cm}\\
(d) Ainda nas notações anteriores, mostre que \(\{X \in \R^{n^2}: XX^t = \id_n\}=:\ ^\infty_{\ d}\mathcal{O}(n)\subset \R^{n^2}\) e calcule \(d\). \\ 
\textcolor{gris}{Segue do \(\blacksquare\) (\emph{Fibras Regulares}) que \(F^{-1}(\{\id_n\}) = \mathcal{O}(n) \) é variadade de clase \(C^\infty\) (\(F\) é \(C^\infty\)) e dimensão 
\[ d = n^2 - \frac{n(n+1)}{2} = \frac{n(n-1)}{2}.\] }\\ 
4. (a) Enunce o \(\blacksquare\) (\emph{Stokes}) \textcolor{gris}{Sejam \(_n\overline{M}\subset \R^n\) compacta, orientada e \(\omega \in \Omega^{n-1}(M)\). Então, tendo \(\partial M\) orientação induzida, 
\[\int_{\overline{M}} d\omega = \int_{\partial M} \omega.\]} \\
(b) Se \(_n\overline{M}\subset \R^m\) é orientável, explique o que significa \(\overline{M}\) ser orientável e como é definida a orientação em \(\partial M\). \\ 
\textcolor{gris}{\(_n\overline{M}\) é \emph{orientavél} sse \(\exists \omega \in \Omega^n(\overline{M}): \forall p \in \overline{M}, \ \omega(p) \neq 0\). \(T_p\partial M \subset T_p\overline{M}\) e \(\dim(T_p \partial M)=n-1\), portanto \(\exists ! \nu(p) \in T_p\overline{M}: \nu(p) \perp T_p\overline{M}\) e aponta pra fora de \(\overline{M}\). Então, \(\partial M\) é orientavél com orientação induzida \(\eta \in \Omega^{n-1}(\partial M)\), 
\[\eta : (v_1, \ldots, v_{n-1}) \mapsto \omega(\nu(p), v_1, \ldots, v_{n-1}).\]}

